%-------------------------
% Resume in Latex -- TikTok Search 2027 Internship
% Target:
% LLM/Multimodal AI Engineer Intern (TikTok Search) - 2027 Start
%
% v19:  Two edits, both requested.
%       (a) Focus line REMOVED from the header block. Availability stays and
%           is now the whole third line, so the "$\;|\;$" separator that used
%           to join them was removed with it. Knock-on effect worth knowing:
%           the v17 note "do not keep the Focus phrase if the intent-taxonomy
%           bullet goes" is now moot, but the reverse dependency is live --
%           nothing in the header points at that bullet any more, so the
%           query-intent taxonomy under Trip.com is the sole carrier of the
%           query-understanding signal. It is doing more work than before.
%           This also frees ~1 line of vertical space at the top of the page;
%           because \raggedbottom top-aligns, that space lands at the BOTTOM,
%           under Leadership. Nothing reflows.
%       (b) Footer added: "Last updated: Aug 2026", bottom-right, \scriptsize
%           italic. Requires a \footskip, since the page had no footer before
%           and \textheight is extended -- see the note at the fancyhdr block.
%           Not compiled/measured. VERIFY once by eye that the footer sits on
%           the paper and does not touch the last Leadership line.
%
% v18:  CONTENT PASS, no spacing changes. Four edits.
%       (a) RETRACTION of v17(c). "(offered)" -> "(reserve tutor)".
%           The CS2103/T email states the opposite of an offer: ~90% of
%           positions went to higher-ranked applicants, the chance of an
%           offer this semester is "rather low", and the coordinator
%           explicitly authorised one phrase and one phrase only --
%           "you are allowed to state on your resume that you were a
%           'reserve tutor'". v17's framing was not optimistic wording,
%           it was contrary to fact, and the accompanying "declined for a
%           scheduling clash" note was a rehearsed line for a question
%           that has a true answer. Both are removed. Do not reinstate.
%           "reserve tutor" is a ROLE, not a status flag -- which is why
%           it reads better than a bare "(reserve)" while staying exact.
%       (b) Availability: "(9 months)" -> "Full-time". The month count is
%           derivable from the dates already on the line; full-time vs.
%           part-time is not, and the JD asks for availability by name.
%           Net -1 char, so the line does not reflow.
%       (c) Trip.com bullet 3: "1,950 evaluations across 5 models" ->
%           the full experiment matrix (questions x models x runs). The
%           lost signal in v17 was "3 runs per model-question pair", i.e.
%           that LLM grading has variance and was measured, not assumed.
%           Written as a product so the reader sees a designed matrix
%           rather than one large number. Costs ~30 chars.
%           >>> VERIFY: 130 is BACK-DERIVED from 1950/5/3. If the real
%           >>> item count differs, fix the first factor. <<<
%       (d) Trip.com bullet 1: "travel agent evaluation system" ->
%           "LLM travel-agent evaluation system". v17 dropped "LLM",
%           which is the JD's own noun ("evaluation modules" for an AI
%           Search Agent); bare "travel agent" can be misread as the
%           human occupation. The hyphen does the disambiguating work.
%       IF THE PAGE OVERFLOWS: (c) and (d) are the additions. Take the
%       space back from Trip.com bullet 1, whose "built to produce
%       rankings that survive re-runs" and "addressing the
%       non-reproducibility of LLM-as-judge grading" say the same thing
%       twice. Do NOT cut the piano/climbing line to pay for this.
%
% v17:  CONTENT PASS, no spacing changes. Four edits, all JD-driven:
%       (a) Trip.com bullet 2: "taxonomy of user information needs" ->
%           "query intent taxonomy". Same claim, the JD's vocabulary
%           ("intent analysis", "query understanding"). Net -2 chars, no
%           reflow. This is the page's only query-understanding evidence.
%       (b) Harmonia bullet 1: named the task (four-way review quality
%           classification on Google location reviews). The pipeline was
%           described but never said what it classified, which left the
%           ~92% in bullet 2 unanchored. Costs ~half a line -- the only
%           edit here with a page-fit risk.
%       (c) Education: "(reserve)" -> "(offered)". "reserve" reads as
%           "not selected"; two TA offers in year 2 is the actual signal.
%       (d) Focus line: added "query understanding", which (a) now
%           supports. Dropped nothing -- "retrieval quality" stays.
%       REMOVED: the commented-out Hybrid Retrieval / reranking slot.
%       Decision taken to NOT ship a ranking experiment: IR/ranking sits
%       in the JD's Preferred list, not Minimum, and a weekend benchmark
%       with untuned baselines interviews worse than an honest gap.
%       If that changes, the block is in v16 of this file.
%
% v3.1: spacing pass (margins, section formatting, macro spacing)
% v4:   content pass -- Education condensed to 2 items ("Awards" label),
%       Trip.com 6 bullets -> 4, bold density cut to one span per bullet.
% v5:   FIX -- \resumeItem now uses %-continuations so a bullet whose last
%       line reaches the right margin no longer emits a phantom blank line.
%       Section lead-in tightened 6pt -> 3pt.
% v6:   Added \titlespacing for \section -- the actual control for the gap
%       between sections. Header block pulled up (-2pt -> -6pt).
% v7:   UNIFORM SECTION GAPS. All lists now have topsep/partopsep pinned to 0,
%       so \titlespacing is the single source of truth for the gap. The Skills
%       block additionally gets \vspace{6pt} because -- unlike every other
%       section -- it does not end with \resumeItemListEnd and so was not
%       receiving that macro's trailing space. Measured gap above each of the
%       five section titles is now equal.
%       TO LOOSEN/TIGHTEN EVERYTHING: change the middle number in \titlespacing.
%       If you do, change the Skills \vspace by the same amount to stay uniform.
% v8:   Body block raised. Space was taken from INSIDE the header block
%       (name \Huge -> \LARGE, the two \vspace{3pt} -> 2pt) rather than from
%       the header->EDUCATION gap, so all five section gaps stay equal.
%       topmargin -0.74in -> -0.79in shifts everything up; the freed space
%       lands under Leadership because \raggedbottom top-aligns the page.
%       Result: top margin 0.21in, bottom margin 0.31in.
% v9:   Name restored to \Huge. The header->EDUCATION gap is now DELIBERATELY
%       smaller than the section-to-section gap (\vspace{-14pt} after the
%       header block). This is not an inconsistency: the header is a light
%       block, the body is dense, and a light/heavy boundary wants less air
%       than a heavy/heavy one. The five SECTION-to-SECTION gaps remain equal
%       -- that is the constraint that actually matters.
%       Knob: the \vspace after \end{center}. -10pt loose, -14pt current,
%       -18pt tight. Every 4pt moved here lands under Leadership.
% v10:  JD keyword pass. (a) "automated scoring" -> "automated answer-quality
%       scoring", matching the JD's "answer quality evaluation" verbatim.
%       (b) Bold moved from engagement-weighted sampling onto the taxonomy
%       phrase -- the taxonomy is the only query-understanding evidence on the
%       page and should be what the eye lands on; count restored to 14.
%       (c) NAACL split out of bullet 1 (which carried four ideas over five
%       lines) into its own bullet, so the first-author signal has a landing
%       spot. Costs one line.
%       (d) Skills: "RAG & vector retrieval" -> "information retrieval (...)",
%       covering the JD's IR/ranking vocabulary without overclaiming.
%       If the page overflows, cut the piano/climbing line in Leadership.
% v12:  (a) Awards + GPA merged onto one line, GPA in the tail position --
%       reclaims the line that pushed the page over, and puts the eye on the
%       scholarship rather than on a mid-range number.
%       (b) SECTION GAP consolidated. It was being fed by three sources at
%       once (resumeItemListEnd's \vspace, the Skills block's \vspace, and
%       titlespacing's "before"), which is why lowering only the last one
%       barely moved it. The first two are now 0 and "before" is 0 as well.
%       ONE KNOB: \titlespacing's middle arg. 0pt = current (tight),
%       4pt = moderate. Do not reintroduce trailing \vspace anywhere.
%       Note the deliberate asymmetry: title->content is 6pt (titlespacing
%       "after" 4pt + \resumeSubheading's 2pt) while content->next-title is
%       ~0pt. A title should bind downward to the block it labels.
% v13:  Title->content gap equalised. Skills and Leadership are bare itemize
%       blocks with no \resumeSubheading, so they were missing that macro's
%       leading \vspace{2pt} and sat 2pt tighter under their titles than
%       Education/Experience/Projects did. Their topsep is now 2pt.
%       If \resumeSubheading's leading \vspace ever changes, change these two
%       topsep values to match, or the asymmetry comes back.
% v16:  MEASURED FROM PDF COORDINATES, not from rasterised pixels.
%       v14/v15 were tuned against pixel ink-gaps, which are corrupted by
%       descenders and antialiasing -- they reported 7.68/9.84/7.44/7.20 for
%       four boundaries that pdftotext -bbox shows are ALL EXACTLY 10.05pt.
%       Two rounds of "fixes" were chasing that noise. The four body sections
%       never needed touching.
%       The one real outlier was Education at 2.08pt, caused by the header
%       pull-up. \vspace{-14pt} -> \vspace{-6pt}. All five now read 10.05pt.
%       The v9 rationale for a deliberately tighter header boundary is hereby
%       withdrawn: uniform reads better here.
%       HOW TO VERIFY (do not trust your eye or a screenshot at 2pt scale):
%         pdflatex resume.tex && python3 bb.py resume.pdf
%       bb.py groups `pdftotext -bbox` words into lines and prints the gap
%       above each section title. Anything within 0.1pt is exact.
%
% v15:  The OTHER half of the boundary. v14 equalised title->content but left
%       content->title alone, which measured:
%         Edu 4.08 | Exp 7.68 | Proj 9.84 | Skills 7.44 | Leadership 3.12
%       Leadership was half of everything else for the same structural reason
%       as v14: every other section ends via \resumeItemListEnd, whose list has
%       topsep=4pt and re-emits it at \endlist. The Skills block is a bare
%       itemize with topsep=0, so it emitted nothing. These two contributions
%       SUM (they are not \addvspace max), hence the exact 4pt shortfall.
%       Fixed with \vspace*{4pt} after the Skills \end{itemize}.
%       Now: Exp 7.68 | Skills 7.44 | Leadership 7.20 -- 0.48pt spread.
%       TWO REMAINING NON-UNIFORMITIES, BOTH INTENTIONAL:
%       - Projects reads 9.84. The line above it ("across revisits.") has no
%         descender; the others end in g/y/p. That 2.40pt IS the descender
%         depth. Baselines are uniform. Do not "fix" this -- padding it would
%         make the baseline grid wrong to correct an artifact of one glyph.
%       - Education reads 4.08 by design: the header is a light block and a
%         light/heavy boundary wants less air (see v9).
%
% v14:  MEASURED, not guessed. Rasterised the PDF at 300dpi and measured the
%       rule-to-first-ink gap under each of the five section titles. v13 read:
%         Education 5.04pt | Experience 4.80 | Projects 3.12 | Skills 1.92 |
%         Leadership 1.92
%       Two separate faults, not one:
%       (1) \resumeProjectHeading lacked the \vspace{2pt} that
%           \resumeSubheading has, so Projects sat 2pt tight. Added it, and
%           took the same 2pt back out of \projGap (6pt -> 4pt) so the gap
%           BETWEEN projects is unchanged.
%       (2) topsep is INERT in the Skills/Leadership blocks -- verified by
%           sweeping it 0pt -> 20pt with no movement at all. Those blocks are
%           written \small{\item{...}}, which wraps \item inside a group and
%           swallows the list's opening skip. Fixed with a rigid \vspace*{3pt}
%           after the section title, which does not depend on list mechanics.
%       Now measures 5.04 / 4.80 / 5.04 / 4.80 / 4.80 -- 0.24pt spread.
%       IF YOU EDIT: topsep will not help in those two blocks. Change the
%       \vspace*{3pt} instead.
%-------------------------

\documentclass[letterpaper,10pt]{article}

\usepackage{latexsym}
\usepackage[empty]{fullpage}
\usepackage{titlesec}
\usepackage{marvosym}
\usepackage[usenames,dvipsnames]{color}
\usepackage{enumitem}
\usepackage[hidelinks]{hyperref}
\usepackage{fancyhdr}
\usepackage[english]{babel}
\usepackage{tabularx}
\input{glyphtounicode}

\pagestyle{fancy}
\fancyhf{}
\fancyfoot{}
\renewcommand{\headrulewidth}{0pt}
\renewcommand{\footrulewidth}{0pt}

% [v19] FOOTER. One line of small text, bottom-right. Deliberately \scriptsize
% + \itshape so it reads as metadata and not as resume content.
% NOTE: \textheight is extended by 1.62in below, so the text block already
% runs close to the bottom edge. \footskip is the ONLY thing keeping the
% footer on the paper -- it is the distance from the bottom of the text block
% down to the footer baseline. Measured bottom margin is ~0.31in (~22pt), so
% 13pt leaves ~9pt of paper under the footer. Do not raise it much past 16pt
% without re-measuring, and if you ever increase \textheight, decrease this.
% TO EDIT THE DATE: change the string below, nothing else.
\fancyfoot[R]{\scriptsize\itshape Last updated: Aug 2026}
\setlength{\footskip}{13pt}

%----------MARGINS----------
\addtolength{\oddsidemargin}{-0.45in}
\addtolength{\evensidemargin}{-0.45in}
\addtolength{\textwidth}{0.8in}
\addtolength{\topmargin}{-0.66in}
\addtolength{\textheight}{1.1in}

\urlstyle{same}
\raggedbottom
\raggedright
\setlength{\tabcolsep}{0in}

% [v19] fancyhdr 在 \pagestyle{fancy} 执行时就把 \textwidth 拷进 \headwidth，
% 而那句在上面 MARGINS 之前，所以页脚宽度停在 fullpage 的原始值、比正文窄 0.8in。
% 必须在 textwidth 改完之后重新同步一次。
\setlength{\headwidth}{\textwidth}

%----------SECTION FORMATTING----------
\titleformat{\section}{
  \scshape\raggedright\large
}{}{0em}{}[\color{black}\titlerule \vspace{-2pt}]

% [v12] SECTION GAP -- SINGLE CONTROL POINT.
% The gap above a section title was never just this number: it was
%   \resumeItemListEnd's trailing \vspace  +  titlespacing's "before"
%   +  itemize's own trailing skip  (partly offset by the -2pt after the rule).
% Those trailing \vspace values are now 0 (see \resumeItemListEnd and the
% Skills block), so the "before" arg below is the ONLY knob. Raise it to
% loosen, lower it to tighten. Args: {indent}{before}{after}.
\titlespacing*{\section}{0pt}{4pt}{4pt}

\pdfgentounicode=1

%----------CUSTOM COMMANDS----------
\newcommand{\resumeItem}[1]{%
  \item\small{{#1}}%
}

\newcommand{\resumeSubheading}[4]{
  \vspace{2pt}\item
    \begin{tabular*}{0.97\textwidth}[t]{l@{\extracolsep{\fill}}r}
      \textbf{#1} & #2 \\
      \textit{\small #3} & \textit{\small #4} \\
    \end{tabular*}\vspace{-3pt}
}

\newcommand{\resumeProjectHeading}[2]{
  \vspace{2pt}\item
    \begin{tabular*}{0.97\textwidth}{l@{\extracolsep{\fill}}r}
      \small #1 & #2 \\
    \end{tabular*}\vspace{-3pt}
}

\renewcommand\labelitemii{$\vcenter{\hbox{\tiny$\bullet$}}$}

\newcommand{\resumeSubHeadingListStart}{
  \begin{itemize}[leftmargin=0.15in, label={}, topsep=0pt, partopsep=0pt,
                  parsep=0pt, itemsep=0pt]
}

\newcommand{\resumeSubHeadingListEnd}{
  \end{itemize}
}

\newcommand{\resumeItemListStart}{
  \begin{itemize}[leftmargin=0.15in, itemsep=1pt, parsep=0pt, topsep=4pt, partopsep=0pt]
}

% [v12] Trailing \vspace zeroed. This fires after the LAST bullet of every
% section too, where it was silently padding the section gap. Section
% separation is now owned solely by \titlespacing.
\newcommand{\resumeItemListEnd}{
  \end{itemize}
}

% [v11] Gaps BETWEEN entries inside a section. Placed by hand between entries
% (never before the first one, never after the last), so section boundaries
% stay tight while entries breathe. Experience entries are longer and denser,
% so they get more air than Projects.
\newcommand{\entryGap}{\vspace{6pt}}
\newcommand{\projGap}{\vspace{3pt}}

%-------------------------------------------
%%%%%% RESUME STARTS HERE %%%%%%%%%%%%%%%%%%
%-------------------------------------------

\begin{document}

%----------HEADING----------
\begin{center}
    \textbf{\Huge \scshape Hongshan Chen} \\ \vspace{2pt}
    \small
    (+65) 9195 4027 $|$
    \href{mailto:e1471281@u.nus.edu}{\underline{e1471281@u.nus.edu}} $|$
    \href{https://github.com/ChenHongshan333}{\underline{GitHub}} $|$
    \href{https://www.linkedin.com/in/hongshan-chen-606928366/}{\underline{LinkedIn}} $|$
    \href{https://chenhongshan333.github.io/}{\underline{Portfolio}}
    \\ \vspace{2pt}
    \small
    % [v19] Focus line removed by request. The separator "$\;|\;$" went with
    % it -- Availability now stands alone on this line and must NOT carry a
    % leading bar. If Focus is ever reinstated, restore the separator too.
    % [v18] "(9 months)" dropped: derivable from the dates. "Full-time" is
    % not derivable, and the JD asks for availability explicitly.
    \textbf{Availability:} Full-time, May 2027 -- Jan 2028
\end{center}

\vspace{-12pt}


%-----------EDUCATION-----------
\section{Education}

\resumeSubHeadingListStart

  \resumeSubheading
    {National University of Singapore}{Aug 2024 -- Jul 2028}
    {Bachelor of Computing in Computer Science, minoring in Interactive Media}{Singapore}

    \resumeItemListStart

      % [v12] Awards + GPA on one line, GPA in the tail position: the
      % scholarship is the stronger signal and should own the eye's landing
      % point. Merging also reclaims the line that pushed the page over.

      
      \resumeItem{
        \textbf{Awards:} NUS Science and Technology Undergraduate Scholarship
      }%

       
      \resumeItem{
        \textbf{GPA:} 4.30/5.00
      }%
      

    \resumeItemListEnd

\resumeSubHeadingListEnd


%-----------EXPERIENCE-----------
\section{Experience}

\resumeSubHeadingListStart

  %----- TRIP.COM -----
  \resumeSubheading
    {Algorithm Intern $|$ LLM Agent Evaluation, Hotel R\&D}
    {Jun 2026 -- Aug 2026}
    {Trip.com Group (Ctrip)}
    {Shanghai, China}

    \resumeItemListStart

      % [v4] old bullets 1 + 2 + 6 merged. The answer-key design is now a
      % clause inside the system description; the NAACL target is a tail phrase.
      % [v18] "LLM" restored and "travel-agent" hyphenated. The JD's noun for
      % this is an agent evaluation module; bare "travel agent" reads as the
      % human occupation. FIRST PLACE TO CUT if the page overflows -- the
      % "survive re-runs" clause and the "non-reproducibility" clause are the
      % same claim stated twice.
      \resumeItem{
        Independently owned an \textbf{end-to-end LLM travel-agent evaluation system} --- corpus
        construction, question generation, answer keys grounded in mechanically decidable
        attributes, automated answer-quality scoring, and leaderboards --- built to produce
        rankings that survive re-runs, addressing the non-reproducibility of LLM-as-judge
        grading.
      }%

      % [v17] "taxonomy of user information needs" -> "query intent taxonomy".
      % Identical claim (information need == intent in the IR literature), but
      % in the JD's own vocabulary: "intent analysis", "query understanding".
      % This is the page's ONLY query-understanding evidence -- it carries the
      % bold, and it is what the Focus line points at.
      \resumeItem{
        Derived a \textbf{query intent taxonomy} (room type,
        restaurants, cancellation policy, \dots) from real user discussions via open coding
        with entropy-based saturation analysis, then applied engagement-weighted sampling to
        align benchmark coverage with the observed query distribution.
      }%

      % [v18] Matrix written out as a product. The point is not the size of
      % 1,950 -- it is the "x 3 runs", which says the variance of LLM grading
      % was measured rather than assumed. A bare total hides that factor.
      % >>> 130 IS BACK-DERIVED (1950 / 5 / 3). Replace with the true item
      % >>> count before sending. A smaller number is fine, arguably better:
      % >>> the tail of this bullet is what justifies it. <<<
      \resumeItem{
        Validated reliability through \textbf{sample-size and ranking-stability analysis}
        (130 questions $\times$ 5 models $\times$ 3 runs = 1,950 evaluations), establishing
        how many items are needed before the leaderboard ordering stops moving.
      }%

      \resumeItem{
        Engineered a \textbf{concurrent Python pipeline} (async workers, Playwright, cascaded
        LLM filtering) that crawls, deduplicates, and validates real user discussions into a
        structured retrieval corpus.
      }%

      \resumeItem{
        Preparing a \textbf{first-author manuscript}, with submission planned for late 2026.
      }

    \resumeItemListEnd

  \entryGap

  %----- NUS RESEARCH -----
  \resumeSubheading
    {Undergraduate Researcher $|$ NUS Odyssey Research Programme
    }
    {May 2026 -- Present}
    {National University of Singapore}
    {Singapore}

    \resumeItemListStart

        \resumeItem{
            Co-designed a \textbf{long-horizon evaluation benchmark for video world models},
            testing whether generated videos preserve evidence and maintain a consistent 3D world across camera departure, occlusion, viewpoint changes, and revisits; manuscript in preparation for submission.
        }%

      \resumeItem{
        Generated controlled synthetic data in \textbf{Unreal Engine 5} (scripted camera
        trajectories, revisit paths, per-object segmentation) --- chosen over in-the-wild video so
        that failures are attributable to the model rather than to uncontrolled scene variation.
      }%

      \resumeItem{
        Contributed to designing \textbf{quantitative metrics and evaluation protocols}
        for reproducibly measuring these failure modes across models.
      }%

    \resumeItemListEnd

\resumeSubHeadingListEnd


%-----------PROJECTS-----------
\section{Projects}

\resumeSubHeadingListStart

  %----- TIKTOK TECHJAM -----
  \resumeProjectHeading
    {
      \textbf{Harmonia} $|$ \textbf{TikTok TechJam 2025}
      $|$
      \emph{FAISS, RAG, CLIP, vLLM}
      \href{https://devpost.com/software/harmonia-fnxts1}
      {\underline{[Project Link]}}
    }
    {Aug 2025 -- Sep 2025}

    \resumeItemListStart

      % [v17] Task named. The pipeline was fully described but never said what
      % it classified, which left bullet 2's ~92% with nothing to attach to.
      % UGC quality classification is also the closest thing on this page to
      % what the target team actually works on -- worth the half line.
      \resumeItem{
        Led architecture of a \textbf{4-stage cascaded multimodal pipeline} for four-way
        review-quality classification (valid vs.\ three violation types) on Google location
        reviews: rule-based filtering $\rightarrow$ FastBERT with entropy-based confidence
        routing $\rightarrow$ Qwen3-8B + CLIP text--image alignment $\rightarrow$ RAG
        verification over a FAISS index, served via vLLM --- routing cheap traffic to cheap
        models and escalating only on low-confidence samples.
      }%

      \resumeItem{
        Profiled the \textbf{accuracy--latency frontier}: $\sim$92\% accuracy at
        $\sim$150ms/sample, roughly 5 points above the latency-matched blend of the two
        endpoint models --- evidence that confidence routing beats naive mixing.
      }

    \resumeItemListEnd

  \projGap

  %----- NETEASE -----
  \resumeProjectHeading
    {
      \textbf{NetEase Cloud Music Support Agent}
      $|$
      \emph{Spring Boot, RAG, Redis}
      \href{https://chenhongshan333.github.io/NetEase-music-agent-backend-demo/}
      {\underline{[Project Link]}}
    }
    {Dec 2025 -- Jan 2026}

    \resumeItemListStart

      \resumeItem{
        Built a \textbf{RAG support-agent backend} in Spring Boot with top-K retrieval and query
        normalisation; added a no-evidence refusal gate that short-circuits generation when
        retrieval returns zero hits, trading answer coverage for groundedness so the agent
        declines rather than fabricates.
      }%

      \resumeItem{
        Wrote a load-testing harness (500 requests @ concurrency 20) that \textbf{diagnosed a JDBC
        connection-pool ceiling}, raising throughput 3.5 $\rightarrow$ 8.3 req/s; added
        retry/circuit-breaker, idempotent writes, and Prometheus metrics.
      }

    \resumeItemListEnd

\resumeSubHeadingListEnd


%-----------TECHNICAL SKILLS-----------
\section{Technical Skills}

\vspace*{3pt}

% [v14] The \vspace*{3pt} above -- NOT topsep -- is what sets this block's
% gap under the title. topsep does nothing here (see header note).
\begin{itemize}[leftmargin=0.15in, label={}, itemsep=2pt,
                topsep=0pt, partopsep=0pt, parsep=0pt]
  \small{\item{
    \textbf{AI \& ML}{: offline evaluation \& metric design, statistical experiment design,
    information retrieval (vector search, RAG, retrieval corpus construction),
    LLM agents (planning, tool use), multimodal learning (CLIP), video understanding,
    prompt engineering} \\[2pt]
    \textbf{Programming}{: Python, C++, Java, C, SQL, JavaScript} \\[2pt]
    \textbf{Tools}{: PyTorch, FAISS, vLLM, NumPy, Spring Boot, Redis, MySQL, Docker,
    Playwright, async Python, Git/GitHub Actions, Linux/CLI, Unreal Engine 5}
  }}
\end{itemize}

\vspace*{4pt}


%-----------LEADERSHIP & ACTIVITIES-----------
\section{Leadership \& Activities}

\vspace*{3pt}

% [v4] kept, but compressed to two lines. If the page overflows, the second
% line (piano / climbing) is the first thing to cut -- the FFC line stays.
% [v14] Gap under the title comes from the \vspace*{3pt} above, not topsep.
\begin{itemize}[leftmargin=0.15in, label={}, itemsep=2pt,
                topsep=0pt, partopsep=0pt, parsep=0pt]

  \small{\item{
    \textbf{Logistics Head, SoC Freshmen Finale Camp 2026} ---
    led logistics planning, vendor coordination, and budget tracking.
  }}

  \small{\item{
    \textbf{NUS Piano Ensemble} --- annual ensemble concerts.
    $|$
    \textbf{NUS Climbing Club} --- SNCS Level 2 Sport Climbing certification.
  }}

  

\end{itemize}

%-------------------------------------------
\end{document}